% Options for packages loaded elsewhere
% Options for packages loaded elsewhere
\PassOptionsToPackage{unicode}{hyperref}
\PassOptionsToPackage{hyphens}{url}
\PassOptionsToPackage{dvipsnames,svgnames,x11names}{xcolor}
%
\documentclass[
  letterpaper,
  DIV=11,
  numbers=noendperiod]{scrartcl}
\usepackage{xcolor}
\usepackage{amsmath,amssymb}
\setcounter{secnumdepth}{5}
\usepackage{iftex}
\ifPDFTeX
  \usepackage[T1]{fontenc}
  \usepackage[utf8]{inputenc}
  \usepackage{textcomp} % provide euro and other symbols
\else % if luatex or xetex
  \usepackage{unicode-math} % this also loads fontspec
  \defaultfontfeatures{Scale=MatchLowercase}
  \defaultfontfeatures[\rmfamily]{Ligatures=TeX,Scale=1}
\fi
\usepackage{lmodern}
\ifPDFTeX\else
  % xetex/luatex font selection
\fi
% Use upquote if available, for straight quotes in verbatim environments
\IfFileExists{upquote.sty}{\usepackage{upquote}}{}
\IfFileExists{microtype.sty}{% use microtype if available
  \usepackage[]{microtype}
  \UseMicrotypeSet[protrusion]{basicmath} % disable protrusion for tt fonts
}{}
\makeatletter
\@ifundefined{KOMAClassName}{% if non-KOMA class
  \IfFileExists{parskip.sty}{%
    \usepackage{parskip}
  }{% else
    \setlength{\parindent}{0pt}
    \setlength{\parskip}{6pt plus 2pt minus 1pt}}
}{% if KOMA class
  \KOMAoptions{parskip=half}}
\makeatother
% Make \paragraph and \subparagraph free-standing
\makeatletter
\ifx\paragraph\undefined\else
  \let\oldparagraph\paragraph
  \renewcommand{\paragraph}{
    \@ifstar
      \xxxParagraphStar
      \xxxParagraphNoStar
  }
  \newcommand{\xxxParagraphStar}[1]{\oldparagraph*{#1}\mbox{}}
  \newcommand{\xxxParagraphNoStar}[1]{\oldparagraph{#1}\mbox{}}
\fi
\ifx\subparagraph\undefined\else
  \let\oldsubparagraph\subparagraph
  \renewcommand{\subparagraph}{
    \@ifstar
      \xxxSubParagraphStar
      \xxxSubParagraphNoStar
  }
  \newcommand{\xxxSubParagraphStar}[1]{\oldsubparagraph*{#1}\mbox{}}
  \newcommand{\xxxSubParagraphNoStar}[1]{\oldsubparagraph{#1}\mbox{}}
\fi
\makeatother


\usepackage{longtable,booktabs,array}
\newcounter{none} % for unnumbered tables
\usepackage{calc} % for calculating minipage widths
% Correct order of tables after \paragraph or \subparagraph
\usepackage{etoolbox}
\makeatletter
\patchcmd\longtable{\par}{\if@noskipsec\mbox{}\fi\par}{}{}
\makeatother
% Allow footnotes in longtable head/foot
\IfFileExists{footnotehyper.sty}{\usepackage{footnotehyper}}{\usepackage{footnote}}
\makesavenoteenv{longtable}
\usepackage{graphicx}
\makeatletter
\newsavebox\pandoc@box
\newcommand*\pandocbounded[1]{% scales image to fit in text height/width
  \sbox\pandoc@box{#1}%
  \Gscale@div\@tempa{\textheight}{\dimexpr\ht\pandoc@box+\dp\pandoc@box\relax}%
  \Gscale@div\@tempb{\linewidth}{\wd\pandoc@box}%
  \ifdim\@tempb\p@<\@tempa\p@\let\@tempa\@tempb\fi% select the smaller of both
  \ifdim\@tempa\p@<\p@\scalebox{\@tempa}{\usebox\pandoc@box}%
  \else\usebox{\pandoc@box}%
  \fi%
}
% Set default figure placement to htbp
\def\fps@figure{htbp}
\makeatother





\setlength{\emergencystretch}{3em} % prevent overfull lines

\providecommand{\tightlist}{%
  \setlength{\itemsep}{0pt}\setlength{\parskip}{0pt}}



 


\KOMAoption{captions}{tableheading}
\usepackage{fontspec}
\newfontfamily\ipaFont{Charis}
\newcommand{\ipa}[1]{{\ipaFont #1}}
\makeatletter
\renewcommand{\tableofcontents}{%
  \section*{Table of Contents}%
  \@starttoc{toc}%
}
\makeatother
\makeatletter
\@ifpackageloaded{caption}{}{\usepackage{caption}}
\AtBeginDocument{%
\ifdefined\contentsname
  \renewcommand*\contentsname{Table of contents}
\else
  \newcommand\contentsname{Table of contents}
\fi
\ifdefined\listfigurename
  \renewcommand*\listfigurename{List of Figures}
\else
  \newcommand\listfigurename{List of Figures}
\fi
\ifdefined\listtablename
  \renewcommand*\listtablename{List of Tables}
\else
  \newcommand\listtablename{List of Tables}
\fi
\ifdefined\figurename
  \renewcommand*\figurename{Figure}
\else
  \newcommand\figurename{Figure}
\fi
\ifdefined\tablename
  \renewcommand*\tablename{Table}
\else
  \newcommand\tablename{Table}
\fi
}
\@ifpackageloaded{float}{}{\usepackage{float}}
\floatstyle{ruled}
\@ifundefined{c@chapter}{\newfloat{codelisting}{h}{lop}}{\newfloat{codelisting}{h}{lop}[chapter]}
\floatname{codelisting}{Listing}
\newcommand*\listoflistings{\listof{codelisting}{List of Listings}}
\makeatother
\makeatletter
\makeatother
\makeatletter
\@ifpackageloaded{caption}{}{\usepackage{caption}}
\@ifpackageloaded{subcaption}{}{\usepackage{subcaption}}
\makeatother
\usepackage{bookmark}
\IfFileExists{xurl.sty}{\usepackage{xurl}}{} % add URL line breaks if available
\urlstyle{same}
\hypersetup{
  colorlinks=true,
  linkcolor={blue},
  filecolor={Maroon},
  citecolor={Blue},
  urlcolor={Blue},
  pdfcreator={LaTeX via pandoc}}


\author{}
\date{}
\begin{document}


\section{L2 Spanish Subject Realization: German L1 Learners vs.~Spanish
L1
Natives}\label{l2-spanish-subject-realization-german-l1-learners-vs.-spanish-l1-natives}

This repository contains the data and R scripts for the study:
\textbf{``The subject realization in L2 Spanish by German L1 speakers: A
corpus study''}

\begin{itemize}
\tightlist
\item
  \textbf{Publication DOI:}
  \url{https://doi.org/10.5565/rev/isogloss.437}
\item
  \textbf{Primary Corpus Source:}
  \href{https://cedel2.learnercorpora.com/}{CEDEL2} (Lozano, 2013)
\end{itemize}

\begin{center}\rule{0.5\linewidth}{0.5pt}\end{center}

\subsection{Repository Structure}\label{repository-structure}

This project uses the OSF Component strategy. To replicate the results,
ensure all files are downloaded into a single working directory.

\subsubsection{Component: Data}\label{component-data}

\begin{itemize}
\tightlist
\item
  \textbf{datal2s.csv}: Subject realization data for 82 German L1
  learners of Spanish across all proficiency levels.
\item
  \textbf{datanatives.csv}: A random sample of 17 native Spanish
  speakers used as a comparative baseline.
\end{itemize}

\subsubsection{Component: Analysis}\label{component-analysis}

\begin{itemize}
\tightlist
\item
  \textbf{demographics.R}: Script for participant metadata and learner
  profile summaries.
\item
  \textbf{descriptive.R}: Script for frequency counts.
\item
  \textbf{mupdarfanalysis.R}: The core script implementing the MuPDARF
  deviation analysis.
\end{itemize}

\begin{center}\rule{0.5\linewidth}{0.5pt}\end{center}

\subsection{Data Description}\label{data-description}

The datasets include sentence-level coding for subject realization in
the \textbf{``Chaplin'' task} (silent film retelling).

\subsubsection{Inclusion Criteria}\label{inclusion-criteria}

To isolate contexts where subject realization is optional/variable, the
following were excluded: * Imperatives * Indefinite ``se'' constructions
* Impersonal sentences (e.g., hay, llueve)

\subsubsection{Key Variables Table}\label{key-variables-table}

{\def\LTcaptype{none} % do not increment counter
\begin{longtable}[]{@{}
  >{\raggedright\arraybackslash}p{(\linewidth - 4\tabcolsep) * \real{0.3377}}
  >{\raggedright\arraybackslash}p{(\linewidth - 4\tabcolsep) * \real{0.3247}}
  >{\raggedright\arraybackslash}p{(\linewidth - 4\tabcolsep) * \real{0.3377}}@{}}
\toprule\noalign{}
\begin{minipage}[b]{\linewidth}\raggedright
Variable
\end{minipage} & \begin{minipage}[b]{\linewidth}\raggedright
Description
\end{minipage} & \begin{minipage}[b]{\linewidth}\raggedright
Values
\end{minipage} \\
\midrule\noalign{}
\endhead
\bottomrule\noalign{}
\endlastfoot
\textbf{Subject} & Binary realization & NS (Null), RS (Realized) \\
\textbf{SubjectType} & Detailed classification & dp, nullsub, pron \\
\textbf{SubjectPosition} & Syntactic position & PRE, POST \\
\textbf{Thetarole} & Semantic role & agent, pacient \\
\textbf{InformationStructure} & Discourse status & GI (Given), NE
(New) \\
\textbf{Proficiency} & Learner level & LA to HC \\
\end{longtable}
}

\begin{center}\rule{0.5\linewidth}{0.5pt}\end{center}

\subsection{Reproducibility
Instructions}\label{reproducibility-instructions}

\subsubsection{1. Requirements}\label{requirements}

You will need \textbf{R} and the following libraries installed: *
partykit * dplyr * ggplot2 * stringr * lattice

\subsubsection{2. Setting the
Environment}\label{setting-the-environment}

Because this repository does not include an \texttt{.Rproj} file, you
must manually set your working directory in R to the folder containing
the downloaded files:

\texttt{setwd("path/to/your/downloaded/files")}

\subsubsection{3. Analysis Procedure:
MuPDARF}\label{analysis-procedure-mupdarf}

The analysis follows the \textbf{MuPDARF} (Multifactorial Prediction and
Deviation Analysis using Random Forests) method: 1. \textbf{Baseline}: A
Random Forest is trained on \texttt{datanatives.csv} to establish native
speaker patterns. 2. \textbf{Prediction}: The model predicts the
expected native-like choices for the German L1 learners in
\texttt{datal2s.csv}. 3. \textbf{Deviation}: Analysis of the statistical
deviation between actual learner choices and model-predicted native
choices.

\begin{center}\rule{0.5\linewidth}{0.5pt}\end{center}

\subsection{Citation}\label{citation}

If you use these scripts or data, please cite the original paper:

\textbf{Kocher, Anna (2025). The subject realization in L2 Spanish by
German L1 speakers: A corpus study. Isogloss. Open Journal of Romance
Linguistics, 11(2), 1--34.}
\url{https://doi.org/10.5565/rev/isogloss.437} And the corpus:

\textbf{Lozano, C. (2013). CEDEL2: Design, compilation and reach of a
corpus of L2 Spanish. Second Language Research, 29(1), 135-147.}




\end{document}
